\chapter{Phương pháp}\label{chapter-3-method}

Chương này trình bày chi tiết về phương pháp đề xuất nhằm đảm bảo tính tái lập (reproducibility). Đối tượng nghiên cứu là test-time SDPO {[}1, §5{]}: một bài toán khó duy nhất được cố định, mô hình tự cải thiện thông qua một số lượng nhỏ các bước huấn luyện ngay tại thời điểm suy luận (test-time training, TTT), và tôi tiến hành đo lường hiệu năng khám phá lời giải của nó. Vì privileged context điều khiển bước cập nhật chính là execution feedback (test case thất bại, lỗi khi thực thi), chế độ này được gọi là execution-guided self-distillation, với mục tiêu là tăng tốc quá trình test-time discovery trong tác vụ sinh mã phức tạp. Trong chế độ này, khóa luận đối chiếu hai cách tổ chức bước cập nhật — student-first (baseline của bài báo gốc {[}1{]}) và teacher-first (một biến thể do người hướng dẫn đề xuất) — đồng thời khảo sát cách reprompt template định dạng (formulate) execution feedback (RO1).

Phạm vi kiểm chứng trên các miền (domain scope) được thiết kế bất đối xứng có chủ đích. Sinh mã là miền thực nghiệm cốt lõi (LiveCodeBench v6 {[}6{]}); trong khi miền toán học chỉ đóng vai trò là một pilot và một điểm tương phản (AIME 2026 {[}10{]}), được sử dụng để đặc trưng hóa ranh giới của cơ chế thay vì để chứng minh tính phổ quát của phương pháp. Sự phân chia đó được giữ nhất quán xuyên suốt: mọi thiết lập trên miền toán đều đi kèm với các giới hạn (caveats) tương ứng về mô hình và reference (§3.5, Chương 6).

Hệ thống ký hiệu trong chương này được đồng nhất với bài báo gốc, Hübotter và cộng sự~{[}1{]}, để người đọc có thể dễ dàng đối chiếu song song.

\begin{center}\rule{0.5\linewidth}{0.5pt}\end{center}

\section{Pipeline test-time SDPO}\label{test-time-sdpo-pipeline}

\subsection{Ký hiệu và thiết lập}\label{notation-and-setting}

Cố định một bài toán \(x\) (đề bài). Hai policy chia sẻ cùng một bộ trọng số \(\theta\):

\begin{itemize}
\tightlist
\item
  student \(\pi_\theta(\cdot \mid x)\), chỉ tiếp nhận đề bài, và
\item
  self-teacher \(\pi_\theta(\cdot \mid x, c)\), chính là mô hình đó nhưng được điều kiện hóa thêm trên một privileged context \(c\).
\end{itemize}

Trong SDPO, \(c\) đóng chính xác vai trò "feedback" \(f\) của bài báo gốc: thông tin hồi cứu (một lỗi khi thực thi, một test case thất bại, hay một lời giải reference) giúp mô hình hồi cứu và xác định được lỗi sai của mình {[}1, §2{]}. Vì teacher và student chỉ khác nhau ở input context, đây là quá trình self-distillation thực thụ: mô hình tự học bằng cách kéo phân bố next-token của nó (khi không có context) về phía phân bố của chính nó khi được điều kiện hóa trên \(c\).

Thiết lập test-time tuân theo đúng Hübotter và cộng sự~{[}1, §5{]}. Với mỗi bài \(x\), mô hình được reset về một base checkpoint rồi chạy \(K\) bước TTT chỉ trên bài đó. Reward là verifiable (kiểm chứng được) — chấm điểm theo thang bậc (graded) cho tác vụ code, và nhị phân cho tác vụ toán. Mục tiêu không phải là tổng quát hóa sang bài khác, mà là khám phá lời giải cho \(x\) càng sớm càng tốt; metric được sử dụng ở đây là discovery@k {[}1, Def. 5.1{]} (§3.6).

\subsection{Vòng lặp chung}\label{the-shared-loop}

Cả hai nhánh chia sẻ chung một bộ khung thuật toán (skeleton) giống hệt nhau. Điểm khác biệt duy nhất nằm ở việc quỹ đạo (trajectory) nào được chọn để distill:

\begin{lstlisting}
load base model + LoRA + optimizer (AdamW, lr = 1e-5)
PRE-eval (student pi_theta(*|x) only, N_eval samples)        # baseline
for step in 1..K:
    feedback        <- build_feedback(student attempt, verifier)   # rich feedback
    {trajectories}  <- generate(...)                               # arm-specific
    {distill_targets} <- select(...)                               # arm-specific (the core difference)
    if distill_targets != {}:
        theta <- theta - lr * grad_theta L_KL(distill_targets)                    # one optimizer step
    log step stats -> W&B
POST-eval (student pi_theta(*|x) only, N_eval samples)       # after TTT
report PRE->POST and the discovery curve
\end{lstlisting}

Hình~\ref{fig:two-arms} phác họa hai nhánh này; chúng chỉ tách nhau đúng tại khối được tô đậm.

\begin{figure}[H]
\centering
\begin{tikzpicture}[font=\footnotesize,>=Stealth,semithick,every node/.style={draw,rounded corners,align=center,minimum height=0.6cm,text width=3.6cm,inner sep=2pt}]
\node(x) at (3,5){Problem $x$};
\node(s) at (0,4){Student $\pi_\theta(\cdot\,|\,x)$};
\node(t) at (6,4){Self-teacher $\pi_\theta(\cdot\,|\,x,c)$};
\node(ys) at (0,2.8){$y_{\text{student}}$ (usually wrong)};
\node(yt) at (6,2.8){$\{y_{t_i}\}\to$ verifier + judge};
\node(g) at (6,1.7){good pool $y_{\text{good}}$};
\node(kls) at (0,1.5){KL distil on $y_{\text{student}}$};
\node(klt) at (6,0.5){KL distil on $y_{\text{good}}$};
\node(u) at (3,-0.6){update $\theta$};
\draw[->](x)--(s);\draw[->](x)--(t);\draw[->](s)--(ys);\draw[->](t)--(yt);\draw[->](yt)--(g);\draw[->](ys)--(kls);\draw[->](g)--(klt);\draw[->](kls)--(u);\draw[->](klt)--(u);
\end{tikzpicture}
\caption{Hai nhánh của test-time SDPO. Student-first và teacher-first chia sẻ toàn bộ vòng lặp và chỉ tách nhau ở việc quỹ đạo nào được distill (khối tô đậm).}
\label{fig:two-arms}
\end{figure}

Toàn bộ sự khác biệt này có thể được tóm gọn như sau: student-first distill chính attempt thất bại của student, còn teacher-first distill một quỹ đạo đúng, đã qua thanh lọc, do teacher sinh ra. Mọi thiết lập khác — hàm mất mát, mô hình, hyperparameter — đều được giữ nguyên vẹn nhằm cô lập duy nhất biến số này.

\begin{center}\rule{0.5\linewidth}{0.5pt}\end{center}

\section{Student-first SDPO (baseline, on-policy)}\label{student-first-sdpo-baseline-on-policy}

Đây là thuật toán gốc của Hübotter và cộng sự~{[}1{]}, hiện thực qua \passthrough{\lstinline!SDPOTrainer!} của TRL {[}5{]} (script \passthrough{\lstinline!07\_discovery\_curve.py!}).

Ở mỗi bước, student sample một rollout \(y \sim \pi_\theta(\cdot \mid x)\) theo kiểu on-policy. Verifier chấm \(y\) và sinh ra feedback \(f\). Self-teacher \(\pi_\theta(\cdot \mid x, f)\) sau đó chấm lại \emph{từng token} của chính \(y\): cho \(f\), phân bố next-token đúng lẽ ra phải là gì tại mỗi vị trí? Student được kéo về phía phân bố hồi tưởng này bằng một KL theo từng token:

\[
\mathcal{L}_{\text{SDPO}}(\theta) \;=\; \sum_{t=1}^{|y|} \mathrm{KL}\,\Big(\pi_\theta(\cdot \mid x, y_{<t}) \;\big\|\; \mathrm{stopgrad}\,\pi_\theta(\cdot \mid x, f, y_{<t})\Big).
\]

\passthrough{\lstinline!stopgrad!} chặn gradient chảy qua teacher, nhờ vậy teacher không sụp vào student và bỏ qua \(f\) {[}1, §2{]}. Gradient và xấp xỉ top-K này được chia sẻ với teacher-first, sẽ được nói kỹ hơn ở §3.3.5.

Có một chế độ thất bại trên bài khó, và chính nó thúc đẩy toàn bộ phần tiếp theo. Vì rollout là on-policy, trên một bài mà bản thân student thuần (chưa được điều kiện hóa) hiếm khi giải được, gần như mọi rollout đều sai — feedback vì thế nghèo nàn, tín hiệu distillation yếu và bất ổn. Đây chính là flat-reward trap: không có rollout đúng nào để học theo cả. Hübotter và cộng sự~{[}1, §5{]} cho thấy SDPO vẫn có thể lặp từ rich feedback ngay cả khi độ chính xác teacher ban đầu dưới 1\%, nhưng họ dùng Qwen3-8B; còn trên mô hình 4B yếu hơn dùng ở đây (§3.5), hành vi kẹt-ở-0 lại thực sự xuất hiện — và đó chính xác là thứ teacher-first nhắm tới giải quyết.

\begin{center}\rule{0.5\linewidth}{0.5pt}\end{center}

\section{Teacher-first SDPO (đề xuất)}\label{teacher-first-sdpo-proposed}

\subsection{Động lực}\label{motivation}

Hai vấn đề của student-first là động lực cho biến thể này.

Vấn đề thứ nhất là flat-reward trap vừa nêu ở §3.2: khi student chưa từng tung ra một attempt đúng, thì không có gì đúng để distill cả.

Vấn đề thứ hai là information leak {[}2{]}. Nếu \(c\) chứa một lời giải reference, teacher sẽ chấm lại theo hướng "càng gần reference càng tốt" — nên qua nhiều bước test-time, student dần hội tụ về việc \emph{sao chép} reference thay vì thực sự học lập luận. Kim và cộng sự~{[}2{]} hình thức hóa hiện tượng này: khi \(I(y; c \mid x)\), tức độ giàu thông tin của context, càng cao, mô hình càng đè nén epistemic verbalization và càng suy giảm.

Thay đổi cốt lõi ở đây là đảo ngược thứ tự. Teacher sinh trước; các quỹ đạo đúng và độc lập được lọc ra; và student được distill về phía chúng. KL được tính trên \(y_{\text{good}}\) (do teacher sinh, đã qua lọc), chứ không phải trên \(y_{\text{student}}\) (một attempt thất bại):

\[
\mathcal{L}_{\text{TF}}(\theta) \;=\; \sum_{y \in \mathcal{G}} \; \sum_{t=1}^{|y|} \mathrm{KL}\,\Big(\pi_\theta(\cdot \mid x, y_{<t}) \;\big\|\; \mathrm{stopgrad}\,\pi_\theta(\cdot \mid x, c, y_{<t})\Big),
\]

trong đó \(\mathcal{G}\) là good pool (§3.3.3). Về mặt khái niệm, teacher-first nằm ở đâu đó giữa SDPO (on-policy) và SFT (off-policy). Nó giống SDFT {[}3{]} ở chỗ distill về phía các generation do chính mô hình (được điều kiện hóa trên context) sinh ra, nhưng context ở đây là feedback kiểu SDPO {[}1{]}, chứ không phải một demonstration cố định.

Vì \passthrough{\lstinline!SDPOTrainer!} hard-code sẵn rollout student-first bên trong, teacher-first phải được hiện thực như một training loop tùy chỉnh riêng (script \passthrough{\lstinline!09\_teacher\_first.py!}), tái sử dụng các helper của \passthrough{\lstinline!07!} (\passthrough{\lstinline!build\_privileged\_context!}, \passthrough{\lstinline!build\_dynamic\_feedback!}, \passthrough{\lstinline!evaluate\_solution!}, và phần đánh giá).

\begin{figure}[H]
\centering
\begin{tikzpicture}[font=\footnotesize,>=Stealth,semithick,every node/.style={draw,rounded corners,align=center,minimum height=0.6cm,text width=3.4cm,inner sep=2pt}]
\node(c) at (0,4){privileged context $c$: feedback + few-shot};
\node(t) at (0,3){self-teacher samples $N$};
\node(v) at (0,2){verifier: correctness};
\node(j) at (0,1){judge: is\_copy / independence};
\node(gp) at (4.8,1){good pool (correct \& independent)};
\node(fs) at (4.8,3){few-shot good\_only / good\_bad};
\node(kl) at (4.8,-0.2){KL distil student $\to y_{\text{good}}$};
\draw[->](c)--(t);\draw[->](t)--(v);\draw[->](v)--(j);\draw[->](j)--(gp);\draw[->](gp)--(fs);\draw[->](fs)--(t);\draw[->](gp)--(kl);
\end{tikzpicture}
\caption{Bước teacher-first. Self-teacher sinh dưới privileged context; một verifier và judge lọc các quỹ đạo vào một good pool; các exemplar good (và tùy chọn bad) lái rollout teacher kế tiếp theo kiểu in-context; student chỉ được distill về phía các quỹ đạo good đã lọc.}
\label{fig:tf-step}
\end{figure}

\subsection{Rollout teacher và privileged context}\label{teacher-rollout-and-privileged-context}

Ở mỗi bước, teacher sample \(N\) quỹ đạo (\passthrough{\lstinline!teacher\_n!}) ở nhiệt độ cao để đảm bảo đa dạng:

\[
\{y_{t_1}, \dots, y_{t_N}\} \sim \pi_\theta(\cdot \mid x, c), \qquad c = \underbrace{f}_{\text{feedback}} \;+\; \underbrace{\text{few-shot exemplars}}_{\text{§3.3.4}}.
\]

\passthrough{\lstinline!privileged\_context!} \(c\) chính là hiện thân cụ thể của "feedback" trong {[}1{]}; cái tên thuộc về codebase của khóa luận, còn khái niệm thuộc về bài báo gốc. Nội dung của \(c\) phụ thuộc vào miền (§3.5). Khối few-shot được chèn ngay trước chỉ dẫn kết ("Respond with ONLY\ldots{}"), với tối đa \passthrough{\lstinline!max\_fewshot = 3!} exemplar để context không bị tràn, và số token được ghi log ở mỗi bước.

\subsection{Verifier và judge tạo good pool}\label{verifier-and-judge-produce-the-good-pool}

Mỗi quỹ đạo \(y_{t_i}\) được gán nhãn qua hai giai đoạn.

\textbf{Giai đoạn 1, verifier (correctness).} Với code (LCBv6), \passthrough{\lstinline!evaluate\_solution!} chạy cả public lẫn private test và trả về tỉ lệ test đạt — một graded reward trong \([0,1]\), cho một gradient dùng được để leo dần (ví dụ 0.21 rồi 1.0). Với toán (AIME), việc khớp đáp án cho một binary reward trong \(\{0,1\}\).

\textbf{Giai đoạn 2, judge (independence).} Giai đoạn này đo xem \(y_{t_i}\) có đơn thuần sao chép reference hay không:

\[
\text{good} := (\text{score} \ge 1.0) \,\wedge\, (\text{independent}), \qquad \text{bad} := (\text{copies the reference}).
\]

Hai cách hiện thực judge được dùng, và sau này được ablate riêng (§4.6):

\begin{longtable}[]{@{}
  >{\raggedright\arraybackslash}p{(\columnwidth - 4\tabcolsep) * \real{0.3333}}
  >{\raggedright\arraybackslash}p{(\columnwidth - 4\tabcolsep) * \real{0.3333}}
  >{\raggedright\arraybackslash}p{(\columnwidth - 4\tabcolsep) * \real{0.3333}}@{}}
\caption{Các hiện thực judge: difflib so với LLM (cơ chế và chi phí).}\\
\toprule\noalign{}
\begin{minipage}[b]{\linewidth}\raggedright
Judge
\end{minipage} & \begin{minipage}[b]{\linewidth}\raggedright
Cơ chế
\end{minipage} & \begin{minipage}[b]{\linewidth}\raggedright
Chi phí
\end{minipage} \\
\midrule\noalign{}
\endfirsthead
\toprule\noalign{}
\begin{minipage}[b]{\linewidth}\raggedright
Judge
\end{minipage} & \begin{minipage}[b]{\linewidth}\raggedright
Cơ chế
\end{minipage} & \begin{minipage}[b]{\linewidth}\raggedright
Chi phí
\end{minipage} \\
\midrule\noalign{}
\endhead
\bottomrule\noalign{}
\endlastfoot
\passthrough{\lstinline!difflib!} & độ tương đồng chuỗi \passthrough{\lstinline!SequenceMatcher!} trên code đã chuẩn hóa (bỏ comment và whitespace); coi là copy nếu tương đồng ≥ ngưỡng (≈0.9) & rẻ, tất định \\
LLM & groq \passthrough{\lstinline!llama-3.3-70b!} (code) hoặc \passthrough{\lstinline!glm-4.5-flash!} (toán), một prompt derive-vs-copy trả về \passthrough{\lstinline!is\_copy!} và \passthrough{\lstinline!reasoning\_quality!} 1--5 & chi phí API, mang tính ngữ nghĩa \\
\end{longtable}

\begin{quote}
\textbf{Một ghi chú về judge.} Khi bật thinking (code), mô hình phát ra một reasoning trace tường minh, nên \passthrough{\lstinline!reasoning\_quality!} chấm chính trace đó thay vì bão hòa ở mức 4--5. Vì vậy LLM judge đóng hai vai cùng lúc: \passthrough{\lstinline!is\_copy!} là cổng độc lập quyết định vào good pool hay không, còn \passthrough{\lstinline!reasoning\_quality!} cung cấp tín hiệu nuôi phép đo epistemic-suppression (§4.7). Việc phát hiện sao chép vẫn là bảo đảm correctness chính của nó.
\end{quote}

Có một hành vi đáng nêu ra ngay từ đầu, vì sau này nó trở thành một giới hạn được đo đạc hẳn hoi. Khi không quỹ đạo nào vượt qua được Giai đoạn 2 — tức mọi ứng viên đều bị gắn cờ là copy — good pool bị để rỗng và bước đó không distill gì cả; một copy bị từ chối không bao giờ bị đưa vào pool một cách khiên cưỡng. Trên một bài mà teacher chỉ luôn sinh ra bản sao, student vì vậy không nhận được tín hiệu học nào hết; điều này được xác nhận trong log toán cho idx8 (Phụ lục D.2), và được nêu ra ở đây để cơ chế được minh bạch.

\subsection{Lái teacher bằng few-shot (good\_only so với good\_bad)}\label{few-shot-teacher-steering-good_only-vs-good_bad}

Các exemplar đã gán nhãn được đưa ngược vào prompt của teacher (in-context, không có gradient trên teacher) nhằm lái teacher hướng tới các generation tốt hơn:

\begin{longtable}[]{@{}
  >{\raggedright\arraybackslash}p{(\columnwidth - 4\tabcolsep) * \real{0.3333}}
  >{\raggedright\arraybackslash}p{(\columnwidth - 4\tabcolsep) * \real{0.3333}}
  >{\raggedright\arraybackslash}p{(\columnwidth - 4\tabcolsep) * \real{0.3333}}@{}}
\caption{Lái teacher bằng few-shot: exemplar good\_only so với good\_bad.}\\
\toprule\noalign{}
\begin{minipage}[b]{\linewidth}\raggedright
\end{minipage} & \begin{minipage}[b]{\linewidth}\raggedright
Phương án 2 --- \passthrough{\lstinline!good\_only!}
\end{minipage} & \begin{minipage}[b]{\linewidth}\raggedright
Phương án 1 --- \passthrough{\lstinline!good\_bad!}
\end{minipage} \\
\midrule\noalign{}
\endfirsthead
\toprule\noalign{}
\begin{minipage}[b]{\linewidth}\raggedright
\end{minipage} & \begin{minipage}[b]{\linewidth}\raggedright
Phương án 2 --- \passthrough{\lstinline!good\_only!}
\end{minipage} & \begin{minipage}[b]{\linewidth}\raggedright
Phương án 1 --- \passthrough{\lstinline!good\_bad!}
\end{minipage} \\
\midrule\noalign{}
\endhead
\bottomrule\noalign{}
\endlastfoot
Cơ chế & few-shot dương (chỉ exemplar good) & few-shot tương phản (good + bad, có gắn nhãn) \\
Độ mạnh lái & vừa phải & mạnh hơn \\
Leak & \textbf{sạch} (good đã qua lọc, nên độc lập với reference) & \textbf{dư} (bad ≈ reference, nên nội dung giống-reference lại tái nhập context) \\
\end{longtable}

Vì \passthrough{\lstinline!bad ≈ reference!}, việc đưa các exemplar bad vào prompt — dù đã gắn nhãn "bad" — vẫn cho teacher thấy lại thứ gần với reference một lần nữa. Ablation giữa Phương án 1 và Phương án 2, vì thế, đồng thời cũng là một ablation leak so với no-leak, một phép kiểm trực tiếp cho chính mối lo của người hướng dẫn. Kết quả trên dữ liệu này là một leak-null (§4.6).

\subsection{Bước distillation: reverse KL top-K căn theo token}\label{the-distillation-step-token-aligned-top-k-reverse-kl}

Đây là lõi tính toán, dùng chung cho cả hai nhánh; chỉ khác nhau ở quỹ đạo target (\(y_{\text{student}}\) so với \(y_{\text{good}}\)).

\textbf{Phân bố theo từng token.} Với logit \(z \in \mathbb{R}^V\) (\(V \approx 150\text{k}\)), \(\pi(v) = \mathrm{softmax}(z)_v\). Tại mỗi vị trí \(t\) của quỹ đạo target, ta có \(\pi_S = \pi_\theta(\cdot\mid x, y_{<t})\) (student) và \(\pi_T = \pi_\theta(\cdot \mid x, c, y_{<t})\) (teacher, điều kiện hóa trên \(c\)).

\textbf{Gradient cốt lõi.} Với teacher được tách gradient, gradient của loss theo một student logit là {[}1; dẫn xuất trong \passthrough{\lstinline!con\_sdpo\_loss\_mechanics!}{]}:

\[
\frac{\partial \mathcal{L}}{\partial z_v^{S}} \;=\; \pi_S(v) - \pi_T(v).
\]

Ở bất kỳ đâu teacher tin một token hơn student (\(\pi_T > \pi_S\)), optimizer sẽ nâng logit đó lên. Target là một phân bố đầy đủ \(V\)-chiều cho mỗi token, chứ không phải một scalar duy nhất cho cả chuỗi như GRPO {[}8{]}. Credit assignment dày hơn hẳn này — nhiều tín hiệu hơn khoảng \(T \cdot K\) lần — chính là nguồn của sample efficiency mà SDPO đạt được {[}1{]}.

\textbf{Hướng KL (\(\alpha\)).} Codebase {[}4{]} tham số hóa divergence qua \(\alpha \in [0,1]\), với mixture \(M = (1-\alpha)\pi_S + \alpha\pi_T\): \(\alpha = 0\) là forward KL (mode-covering, giữ lại các epistemic token), \(\alpha = 1\) là reverse KL (mode-seeking, dễ bị đè nén hơn), và \(\alpha = 0.5\) là JSD. Khóa luận dùng \(\alpha = 1.0\) (reverse KL, top-K = 20) để khớp với cấu hình LCBv6 của {[}4{]}. Liệu \(\alpha\) có thực sự điều biến sự đè nén hay không vẫn là một câu hỏi RO2 còn bỏ ngỏ (Chương 6).

\textbf{Xấp xỉ top-K.} Tính KL trên toàn vocabulary quá nặng bộ nhớ ở \(V \approx 150\text{k}\), nên ta chỉ dùng top-20 logit của teacher cộng thêm một tail bucket; student được chiếu lên đúng K chỉ số đó, và KL được lấy trên phân bố \((K{+}1)\)-chiều thu được {[}1, §2.2{]}.

\textbf{Importance sampling.} Vì các sample được rút dưới \(\theta_{\text{old}}\) nhưng loss lại được tính dưới \(\theta\) hiện tại, mỗi per-token loss được scale bởi một ratio đã clip \(r_t = \min(\pi_\theta/\pi_{\theta_\text{old}},\, c_{\text{clip}})\) với \(c_{\text{clip}} = 2.0\).

\textbf{Căn token — điểm dễ xảy ra lỗi nhất.} Prompt student (\passthrough{\lstinline!x + y\_good!}) và prompt teacher (\passthrough{\lstinline!x + c + y\_good!}) có độ dài prefix khác nhau, nên các logit tại các vị trí của \(y_{\text{good}}\) phải được trích riêng ở mỗi bên trước khi lấy KL. Chỉ cần lệch một token ở đây là KL sẽ bị sai lệch một cách âm thầm, nên code luôn assert rằng hai độ dài khớp nhau, và log lại \passthrough{\lstinline!(student\_prefix\_len, teacher\_prefix\_len)!} mỗi lần. Sanity check đã cho kết quả đạt: căn token đúng (\(L\) khớp ở cả hai bên dù prefix khác nhau) và token-mean KL hữu hạn, ở mức hợp lý.

Hyperparameter: AdamW \passthrough{\lstinline!lr = 1e-5!}, \passthrough{\lstinline!kl\_topk = 20!}, \passthrough{\lstinline!kl\_alpha = 1.0!}, \passthrough{\lstinline!is\_clip = 2.0!}. Teacher luôn được tách gradient (self-distillation với trọng số chia sẻ, nên teacher thực chất là forward pass của student dưới context \(c\), chạy dưới \passthrough{\lstinline!no\_grad!}).

\begin{center}\rule{0.5\linewidth}{0.5pt}\end{center}

\section{Reprompt template (T1--T7)}\label{reprompt-templates-t1t7}

Reprompt template quyết định nội dung của \(c\) mà teacher nhìn thấy — thứ trực tiếp thay đổi \(\pi_T\), và qua đó thay đổi cả hướng của gradient distillation. Đây là biến trung tâm của RO1. Để lựa chọn có cơ sở đứng vững (câu hỏi hiển nhiên mà một reviewer sẽ hỏi là "vì sao đúng bảy cái này?"), ta không biện minh từng template một cách riêng lẻ, mà tổ chức không gian thiết kế theo ba chiều dựa trên các công trình trước {[}1, 2{]}, rồi lấy mẫu T1--T7 trải đều trên ba chiều đó.

\begin{longtable}[]{@{}
  >{\raggedright\arraybackslash}p{(\columnwidth - 6\tabcolsep) * \real{0.2500}}
  >{\raggedright\arraybackslash}p{(\columnwidth - 6\tabcolsep) * \real{0.2500}}
  >{\raggedright\arraybackslash}p{(\columnwidth - 6\tabcolsep) * \real{0.2500}}
  >{\raggedright\arraybackslash}p{(\columnwidth - 6\tabcolsep) * \real{0.2500}}@{}}
\caption{Phân loại reprompt-template: T1--T7 trên ba chiều thiết kế.}\\
\toprule\noalign{}
\begin{minipage}[b]{\linewidth}\raggedright
Template
\end{minipage} & \begin{minipage}[b]{\linewidth}\raggedright
Tên
\end{minipage} & \begin{minipage}[b]{\linewidth}\raggedright
Chiều
\end{minipage} & \begin{minipage}[b]{\linewidth}\raggedright
Khác T2 ở
\end{minipage} \\
\midrule\noalign{}
\endfirsthead
\toprule\noalign{}
\begin{minipage}[b]{\linewidth}\raggedright
Template
\end{minipage} & \begin{minipage}[b]{\linewidth}\raggedright
Tên
\end{minipage} & \begin{minipage}[b]{\linewidth}\raggedright
Chiều
\end{minipage} & \begin{minipage}[b]{\linewidth}\raggedright
Khác T2 ở
\end{minipage} \\
\midrule\noalign{}
\endhead
\bottomrule\noalign{}
\endlastfoot
T1 & Tối giản & Lượng thông tin (thấp) & bỏ cấu trúc, chỉ ``incorrect, try again'' \\
\textbf{T2} & \textbf{Chuẩn (anchor)} & \textbf{baseline} & \textbf{failing test + expected + actual + error\_type} \\
T3 & Dài dòng & Lượng thông tin (cao) & T2 + full stack trace + code trước \\
T4 & JSON có cấu trúc & Instruction framing (định dạng) & cùng nội dung T2, mã hóa JSON \\
T5 & Kích thích lập luận & Instruction framing (chẩn đoán) & T2 + ``First, identify root cause, then fix.'' \\
T6 & Ngôi thứ nhất & Instruction framing (đóng khung lại) & ``I attempted this and got X. Let me reconsider.'' \\
T7 & Lịch sử tích lũy & Độ sâu bộ nhớ & tất cả N attempt trước, không chỉ cái mới nhất \\
\end{longtable}

\begin{figure}[H]
\centering
\begin{tikzpicture}[font=\footnotesize,>=Stealth,semithick,every node/.style={draw,rounded corners,align=center,minimum height=0.55cm,inner sep=3pt}]
\node(a) at (0,3){T2 standard (anchor)};
\node(d1) at (-4,1.8){Dim 1: information};
\node(d2) at (0,1.8){Dim 2: framing};
\node(d3) at (4,1.8){Dim 3: memory};
\node(t1) at (-5,0.6){T1};\node(t3) at (-3,0.6){T3};
\node(t4) at (-1.4,0.6){T4};\node(t5) at (0,0.6){T5};\node(t6) at (1.4,0.6){T6};
\node(t7) at (4,0.6){T7};
\draw[->](a)--(d1);\draw[->](a)--(d2);\draw[->](a)--(d3);
\draw[->](d1)--(t1);\draw[->](d1)--(t3);\draw[->](d2)--(t4);\draw[->](d2)--(t5);\draw[->](d2)--(t6);\draw[->](d3)--(t7);
\end{tikzpicture}
\caption{Không gian thiết kế reprompt-template. Mỗi template thay đổi một chiều so với anchor T2: lượng thông tin (T1, T3), instruction framing (T4, T5, T6), hoặc độ sâu bộ nhớ (T7).}
\label{fig:template-space}
\end{figure}

Ba chiều đó, và cơ sở của từng chiều:

\begin{enumerate}
\def\labelenumi{\arabic{enumi}.}
\tightlist
\item
  \textbf{Lượng thông tin} (T1 → T2 → T3). Kim và cộng sự~{[}2{]} cho thấy \(I(y; c \mid x)\) chi phối độ lớn của sự đè nén; khóa luận thay đổi lượng thông tin trong execution feedback dựa trên đó.
\item
  \textbf{Instruction framing} (T4, T5, T6). Hübotter và cộng sự~{[}1, §4.6{]} kết luận rằng các biến thể cú pháp nhỏ không quan trọng, nhưng biến thể ngữ nghĩa hay chỉ dẫn thì chưa từng được kiểm; {[}1, §7{]} nêu chính điều này là future work, gần như nguyên văn ("how \ldots{} the reprompt template \ldots{} influences behavior").
\item
  \textbf{Độ sâu bộ nhớ} (T7). Thiết lập ở §5 của bài báo gốc dùng một cửa sổ trượt cố định chỉ một attempt và không ablate độ sâu; trong khi Kim và cộng sự~{[}2{]} lại cho thấy sự đè nén tích lũy qua các bước, nên lịch sử có thể khuếch đại hoặc làm loãng hiệu ứng này.
\end{enumerate}

\textbf{T2 đóng vai anchor}: mọi template khác chỉ thay đổi đúng một chiều và giữ nguyên phần còn lại, cho một ablation sạch sẽ. Trong bảy template, bốn cái được hiện thực (T1, T2, T5, T6; chuỗi nguyên văn ở Phụ lục B), còn ba cái (T3, T4, T7) vẫn chỉ là các điểm khái niệm trong phân loại. Với compute có sẵn, RO1 chỉ dò được T1/T2/T5 (§4.5); T6 tuy đã hiện thực nhưng chưa dò tới, và T3/T4/T7 được để lại cho hướng phát triển sau (Chương 6). Hai caveat cần lưu ý: T4 (JSON) cũng đồng thời thay đổi mật độ thông tin (một sự chồng lấn giữa chiều 1 và 2), và T7 làm tăng độ dài context nên confound với chi phí compute — cần được báo cáo riêng.

\begin{center}\rule{0.5\linewidth}{0.5pt}\end{center}

\section{Các miền: code (cốt lõi) và toán (pilot)}\label{domains-code-core-and-math-pilot}

Hai miền khác nhau ở reference mode — tức thứ mà teacher coi là "đáp án" — và ở nội dung của privileged context. Như Chương 5 sẽ lập luận, đây chính là biến quyết định teacher-first thành công hay không.

\begin{longtable}[]{@{}
  >{\raggedright\arraybackslash}p{(\columnwidth - 4\tabcolsep) * \real{0.3333}}
  >{\raggedright\arraybackslash}p{(\columnwidth - 4\tabcolsep) * \real{0.3333}}
  >{\raggedright\arraybackslash}p{(\columnwidth - 4\tabcolsep) * \real{0.3333}}@{}}
\caption{So sánh miền: reference và context của code (LCBv6) so với toán (AIME).}\\
\toprule\noalign{}
\begin{minipage}[b]{\linewidth}\raggedright
\end{minipage} & \begin{minipage}[b]{\linewidth}\raggedright
\textbf{Code --- LCBv6 {[}6{]}} (cốt lõi)
\end{minipage} & \begin{minipage}[b]{\linewidth}\raggedright
\textbf{Toán --- AIME 2026 {[}10{]}} (pilot)
\end{minipage} \\
\midrule\noalign{}
\endfirsthead
\toprule\noalign{}
\begin{minipage}[b]{\linewidth}\raggedright
\end{minipage} & \begin{minipage}[b]{\linewidth}\raggedright
\textbf{Code --- LCBv6 {[}6{]}} (cốt lõi)
\end{minipage} & \begin{minipage}[b]{\linewidth}\raggedright
\textbf{Toán --- AIME 2026 {[}10{]}} (pilot)
\end{minipage} \\
\midrule\noalign{}
\endhead
\bottomrule\noalign{}
\endlastfoot
Mô hình & Qwen3-4B {[}7{]}, LoRA r=32, thinking \textbf{ON} & Qwen3-4B {[}7{]}, LoRA r=32, thinking \textbf{ON} \\
Verifier & public/private test → \textbf{graded} \([0,1]\) & answer match → \textbf{binary} \(\{0,1\}\) \\
\passthrough{\lstinline!reference\_mode!} & \passthrough{\lstinline!best\_in\_batch!} (quỹ đạo điểm cao nhất làm proxy) & \passthrough{\lstinline!ground\_truth!} (teacher luôn thấy đáp án đúng) \\
Leak regime & \textbf{null} (reference là best-in-batch + public test, không phải một lời giải reference) & \textbf{cực đoan} (feedback là ``the correct answer is \{gt\}'') \\
\passthrough{\lstinline!privileged\_context!} & execution feedback (test thất bại, expected/actual, loại lỗi) + quỹ đạo best-in-batch đúng & đáp án đúng (một con số); teacher chỉ cần chép \passthrough{\lstinline!\\boxed\{\}!} \\
Reference chứa một phương pháp? & \textbf{Có} (best-in-batch là một \emph{chương trình} = một thuật toán, tổng quát hóa sang input mới) & \textbf{Không} (chỉ một \emph{giá trị}, không phải một suy dẫn) \\
\end{longtable}

Một ghi chú về lựa chọn \passthrough{\lstinline!reference\_text!} cho code: LCBv6 không có trường lời-giải-reference đáng tin cậy, nên khóa luận dùng best\_in\_batch (quỹ đạo điểm cao nhất trong batch của teacher ở mỗi bước) làm reference để judge đo tính độc lập. Điều này kéo theo một tác dụng phụ đáng chú ý: vì reference không phải một lời giải của tác giả mà là output đúng của chính mô hình, đã được kiểm chứng bởi public test, nên code không hề có leak — điều này giải thích vì sao ablation good\_bad cho kết quả leak-null trên code (§4.6), trong khi leak lại đo được rõ ràng trên toán (§4.9).

Ngân sách TTT cũng khác nhau giữa hai miền (giá trị đầy đủ ở Phụ lục A): code chạy 15 bước, eval 16 sample, teacher\_n = 10; toán chạy 3 bước, eval 4 sample, teacher\_n = 4, max\_new\_tokens = 16384 (8192 từng cắt cụt đáp án \passthrough{\lstinline!\\boxed!} và gây ra các điểm 0 sai lệch). Khoảng cách ngân sách này là một biến cần lưu ý khi so sánh hai miền, bên cạnh khác biệt về loại reference.

Vì cả hai miền đều dùng cùng một mô hình Qwen3-4B, một kết quả no-escape trên toán không thể quy cho việc mô hình là một reasoner yếu hơn; loại reference vẫn là khác biệt vận hành chính, dù khoảng cách ngân sách cũng góp phần. Do đó, một caveat toán luôn đi kèm mọi kết quả toán:

\begin{enumerate}
\def\labelenumi{\arabic{enumi}.}
\tightlist
\item
  \textbf{Reference chỉ-có-đáp-số.} AIME chỉ cung cấp một đáp án số, không có lời giải kèm theo, nên privileged context được hard-code để trả về thẳng đáp án, và teacher lúc này chỉ có thể chép lại. MATH-500 (có sẵn một trường \passthrough{\lstinline!solution!}) là con đường để có được một reference toán mang-phương-pháp (Chương 6), nhưng nằm ngoài phạm vi các kết quả hiện tại.
\end{enumerate}

Việc chọn bài dựa trên model pass-rate \(\in (0,1)\) — một frontier do chính mô hình định nghĩa, chứ không phải nhãn độ khó của kỳ thi — vì binary math reward chỉ có phương sai khi mô hình giải được một bài trong khoảng 30--70\% số lần thử (§3.6, §4).

\begin{center}\rule{0.5\linewidth}{0.5pt}\end{center}

\section{Metric và thiết kế so sánh}\label{metrics-and-comparison-design}

\textbf{pass@k / pass\_rate.} Ta báo cáo một ước lượng thực nghiệm của pass@k {[}9{]}: với mỗi điều kiện, ta rút \passthrough{\lstinline!N\_eval!} sample (16 cho code, 4 cho toán) và báo cáo \passthrough{\lstinline!pass\_rate!}, tức tỉ lệ đúng. Cả PRE (trước TTT) lẫn POST (sau TTT) đều được báo cáo, chỉ đánh giá student, \(\pi_\theta(\cdot\mid x)\), không có context nào tại thời điểm đánh giá (nếu không, context sẽ leak thẳng vào metric). Một sample \passthrough{\lstinline!greedy!} tất định duy nhất cũng được báo cáo kèm theo, như một tín hiệu phụ, dù có phần nhiễu.

\textbf{Đường cong discovery.} Theo discovery@k {[}1, Def. 5.1{]} đã nêu ở §3.1: xác suất giải được trong \(k\) attempt đầu, \(P(T \le k)\) với \(T = \min\{k : r(y_k) = 1\}\). Vì generation ở test-time là tuần tự (có chia sẻ trạng thái, có cập nhật trọng số), pass@k kiểu i.i.d. không còn áp dụng được nữa, và discovery@k mới là metric đúng. Ta log lại batch pass-rate ở mỗi bước như một đường cong discovery ở dạng thô.

\textbf{Escape-zero (bằng chứng cho cơ chế).} Định nghĩa vận hành: một seed mà ở đó student-first kẹt cứng ở pass = 0 (hoặc rớt trở lại 0), trong khi teacher-first thoát được khỏi 0. Đây chính là chữ ký trực tiếp của flat-reward trap (§3.2): nhánh on-policy không có rollout đúng nào để học, trong khi teacher nghiêng về off-policy có thể tiêm thẳng một lời giải đúng vào. Việc tái lập được escape-zero là đóng góp mang tính cơ chế quan trọng nhất (§4.3).

\textbf{So sánh matched.} Hai nhánh chạy trên \emph{cùng} một base và \emph{cùng} một seed, nên PRE khớp nhau theo từng seed — cho ta một so sánh ghép cặp (matched), loại bỏ được phương sai ở mức base và mức seed. Trên mỗi cặp matched, ta đếm số win, tie và loss (TF so với SF).

\begin{quote}
\textbf{Một ghi chú về lực thống kê.} Nhánh code được đánh giá ở quy mô trung bình (8 bài × 6 seed), đủ sức cho một Wilcoxon signed-rank test trên các POST pass-rate ghép cặp, thay vì chỉ một sign test thô; kết quả chính được báo cáo với thống kê có lực đúng nghĩa này (§4.3, Chương 6). Miền toán thì vẫn chỉ là một pilot n-nhỏ (n = 2 bài), nơi mọi tuyên bố được giữ ở mức \emph{mô tả và định hướng} ("sơ bộ / ưu thế yếu", chứ không bao giờ "ta chứng minh"), với mỗi phát biểu luôn kèm một caveat về n.
\end{quote}
